% Section 7
\section{The R package}\label{sec:package}

The methods are implemented in the R package \pkg{eiPartialID}, version 0.3.1,
which extends the $2\times2$ package of JKST. The original functions
(\code{generateBounds()}, \code{evaluateBounds()}) are unchanged; with a range
offset of $10^{-5}$ and HC1 variances, the new code reproduces them on JKST's
census example to $10^{-10}$. The LPs are solved with \pkg{lpSolve} and the
robust covariances come from \pkg{sandwich}.

\subsection{The main function}

\code{eiBounds()} computes bounds and confidence intervals under one rung. Its
main arguments and defaults are:
\begin{itemize}
\item the data, as a formula or as \code{x}, \code{t} and unit sizes \code{n}, with
an optional \code{group} (the districts);
\item \code{model}: \code{"M0"} (default), \code{"M1"}, \code{"M1delta"} (with
\code{delta}), \code{"M2"}, \code{"L4"} (with \code{row\_constant}), \code{"L5"},
\code{"L6"};
\item \code{targets}: linear targets as weight matrices (by default every cell rate);
\item \code{range}: \code{"hull"} (default; the group-wise sample hull,
Proposition~\ref{prop:consistency}(c)), \code{"trim"} (dropped units enter
through Duncan--Davis, Proposition~\ref{prop:trim}) or \code{"extrapolate"} (pure
units or a user matrix, with a warning that validity there is assumed,
Proposition~\ref{prop:simplex});
\item \code{weights}: \code{"none"} or \code{"n"} for the regression;
\code{vcov}: \code{"HC3"} (default), \code{"HC1"}, \code{"HC0"};
\item \code{ci}: \code{"clr"} (default; Conjecture~\ref{conj:clr}, with
\code{ndraw = 100} simulated draws), \code{"conservative"}
(Theorem~\ref{thm:ci}(a)) or \code{"none"}; \code{alpha = 0.05};
\item \code{fpc}: the residual allowance for the actual rate,
\code{"bhatia-davis"} (default; Proposition~\ref{prop:bd}) or \code{"popoviciu"}
($S_1$);
\item \code{restrictions}: optional rows (\code{zeros}, \code{order}, \code{sign},
\code{slope}, \code{linear}, \code{external}), and the sign module's \code{assoc}
(signs on $\delta_B$, one-vs-rest or pairwise, and on $\delta_W$, W1 or W2) and
\code{cr} (contextual reinforcement);
\item \code{union}: \code{"pieces"} (print every piece of a union) or \code{"hull"};
\item \code{eta}: the sensitivity parameter of Proposition~\ref{prop:S1};
\code{jkst = TRUE} adds JKST's $\CI_x$ with its revert-to-Duncan--Davis rule;
\code{tests = TRUE} computes the tests of Proposition~\ref{prop:tests}.
\end{itemize}

\paragraph{What a fit reports.} For each target, \code{eiBounds()} returns:
\code{estimate}, the plug-in estimate of the outer identified set for the
model-implied rate (not a confidence interval); \code{ci\_model}, the CLR (or
conservative) interval for $\Bt$; \code{ci\_actual}, the same for $B$ with the
residual allowance; \code{proj\_model} and \code{proj\_actual}, the projection
intervals of Theorem~\ref{thm:ci}(b2)(ii) computed from the same candidates, with
$q=\dim\beta$ and, for $B$, the Bhatia--Davis allowance taken from a projection
interval for $\bar m$ on the rung alone (Proposition~\ref{prop:bd}); and the
Duncan--Davis interval. Intervals for $B$ are intersected with Duncan--Davis. The
rank of $\hat\Sigma$ is printed as a diagnostic; candidates whose gradients load on
a numerically null direction of $\hat\Sigma$ are dropped as a labelled safeguard,
which can only widen the projection interval. A CLR candidate with zero estimated
variance is treated as exact only if its gradient is zero.

With optional restrictions, the rung alone is also fitted and printed beside the
restricted result; with several restrictions it is labelled as their joint effect.
The \code{tests} component holds the emptiness test, the curvature test (M1 against
M2) and the pooling test, each with $\chi^2$, F and restricted-wild-bootstrap
p-values ($B=399$). An empty estimate is reported as a sample diagnostic, not a
test; the emptiness test is the test.

\paragraph{What \code{print()} says about each fit.} Every fit prints a status line
that names the result covering its intervals:
\begin{itemize}
\item \emph{Pooled M0:} the projection intervals carry an asymptotic guarantee
(Theorem~\ref{thm:ci}(b2)(ii)), or the conservative interval does
(Theorem~\ref{thm:ci}(a)) when \code{ci = "conservative"}; with sign or CR rows the
line adds that the guarantee is Theorem~\ref{thm:R10} for the pooled model.
\item \emph{Grouped rungs without sign rows:} an asymptotic guarantee only with few,
large groups (regime R1); with many small groups, a simulation finding.
\item \emph{Grouped rungs with sign or CR rows:} no coverage theorem.
\end{itemize}
CLR lines are labelled ``algorithmic, coverage supported by simulation, not
proved''; for $B$ the label adds that the Bhatia--Davis bound also uses a CLR
interval; with CR, the CLR interval is the union of the CLR branch intervals.
Every sign or CR restriction prints a line marking the module as experimental,
with its proof status. Grouped fits with sign rows print that each sign is one
pooled restriction across groups and that the bound is an outer set whose
sharpness under grouped rungs is not proved; pooled fits may label the estimate
sharp for the regression (Theorem~\ref{thm:R3}(b)). Notes flag near-boundary fits
(Slater slack below 0.005), few units per group (median below 30, where CLR
coverage rests on simulation), unbounded estimates (a hull that is not
full-dimensional), a Bhatia--Davis input interval for $\bar m$ inside $[0,.05]$
or $[.95,1]$ (where the condition of Proposition~\ref{prop:bd} is doubtful and
\code{fpc = "popoviciu"} is recommended), and joint rejections, naming the rejected
combination and saying whether the rung alone is rejected.

\subsection{Other functions}

\begin{itemize}
\item \code{eiRedundancy()} reports what each restriction buys, using the face test
of Proposition~\ref{prop:R2} at the point estimate (never solver duals): the
verdict and the width gain when the restriction is added first (to the rung alone)
and last (to the full stack); the associations already settled by the rung
alone, with their intervals; the CR branches and the connectivity certificate of
Theorem~\ref{thm:R3}(c); EGGHM's descriptive quantities ($\gamma$, $\ER$, $\NM$),
with a warning never to intersect the bracket with a confidence interval; and the
weight concentration of each signed association (the largest unit weight and the
effective number of units).
\item \code{ladder()} fits the rungs from the most to the least plausible
(Theorem~\ref{thm:ladder}(c)).
\item \code{sensitivity(dial = "eta")} and \code{breakdown()} implement
Proposition~\ref{prop:S1}; \code{sensitivity(dial = "slack")} relaxes the sign
beliefs by a slack in probability units, with a breakdown value along a stated
scalar path (Proposition~\ref{prop:N1}(d)).
\item \code{eiModules()} lists the restriction types with their proof status, and
\code{eiBenchmark()} computes coverage on data with known truth.
\end{itemize}

\paragraph{Not implemented.} Sign rows on realised tables (N1), the monotone rung
N2, local contextual rows, the series bound that would tighten the outer set
towards $\ThX$, the finite-sample interval of Theorem~\ref{thm:finite} (because it is
uninformative), covariates and integer tables.
